\chapter{GIỚI THIỆU TỔNG QUAN}
\label{chap:intro}

\section{Bối cảnh chung}
Trong bối cảnh nền công nghiệp 4.0 và sự phát triển mạnh mẽ của các hệ thống sản xuất thông minh, robot công nghiệp đang ngày càng đóng vai trò thiết yếu. Tuy nhiên, thay vì tự động hóa hoàn toàn, xu hướng hiện nay tập trung vào Tương tác Người - Máy (Human-Robot Interaction - HRI) và Tương tác cộng tác (Human-Robot Collaboration - HRC) \cite{zhang2023research}. Trong các môi trường làm việc độc hại, phức tạp hoặc không có cấu trúc rõ ràng (như không gian vũ trụ, lò phản ứng hạt nhân, y tế và tháo dỡ bom mìn), trí tuệ nhân tạo hiện tại chưa thể hoàn toàn thay thế khả năng xử lý tình huống linh hoạt của con người \cite{hetrick2020comparing, luu2025enhancing}. Do đó, việc vận hành và điều khiển robot từ xa (teleoperation) kết hợp với trí thông minh thích ứng của con người là một giải pháp tối ưu, cho phép robot thực hiện các nhiệm vụ phức tạp mà không đòi hỏi con người phải có mặt trực tiếp tại hiện trường.

\section{Điều khiển từ xa qua Thực tế ảo (VR Teleoperation) \& Động lực nghiên cứu}
Các hệ thống điều khiển từ xa truyền thống thường phụ thuộc vào màn hình 2D, bàn phím, hoặc tay cầm điều khiển (joystick). Những giao diện này đòi hỏi người vận hành phải chuyển đổi dữ liệu hình ảnh 2D thành các hành động trong không gian 3D, làm tăng đáng kể tải nhận thức (cognitive load) và dễ gây ra sai sót, dẫn đến hiệu suất công việc giảm sút \cite{hetrick2020comparing}. 

Thực tế ảo (Virtual Reality - VR) đã nổi lên như một công nghệ đột phá giúp "xóa nhòa khoảng cách" vật lý. Công nghệ VR cung cấp môi trường nhập vai 3D, khôi phục nhận thức chiều sâu và nhận thức không gian (spatial awareness) một cách tự nhiên. Thông qua các thiết bị HMD (Head-Mounted Display) và bộ điều khiển theo dõi chuyển động (motion controllers), VR cho phép người dùng vận hành robot bằng các cử chỉ trực quan thay vì các thao tác điều khiển cơ học phức tạp \cite{chen2023comparing, torrejón2026teleoperation}. Động lực chính của luận văn này là khai thác tối đa sức mạnh của môi trường VR nhập vai kết hợp với công nghệ Digital Twin (Bản sao số) để tái tạo trải nghiệm hiện diện (telepresence) với độ chân thực cao nhất, qua đó cải thiện tốc độ, độ chính xác và giảm sự mệt mỏi cho người điều khiển.

\section{Các thách thức hiện tại}
Mặc dù tiềm năng của VR trong điều khiển robot từ xa là rất lớn, các hệ thống hiện tại vẫn đang đối mặt với nhiều rào cản về cả phần cứng lẫn nền tảng phần mềm:
\begin{itemize}
    \item \textbf{Độ trễ và đồng bộ hóa (Latency \& Synchronization):} Độ trễ trong truyền thông mạng và sự thiếu đồng bộ giữa mô hình ảo và robot thực tế là rào cản lớn nhất. Sự chậm trễ giữa thao tác tay của người dùng và phản hồi hình ảnh gây ra hiện tượng chóng mặt (cybersickness) và giảm độ chính xác \cite{luu2025enhancing}.
    \item \textbf{Nút thắt cổ chai của Middleware truyền thống:} Rất nhiều nghiên cứu vẫn đang sử dụng Robot Operating System phiên bản 1 (ROS1) với kiến trúc Master/Slave. Kiến trúc này không đảm bảo tốt tính thời gian thực khi phải truyền tải khối lượng dữ liệu khổng lồ (như hình ảnh camera, point cloud, và trạng thái khớp) liên tục giữa môi trường VR và robot \cite{zhang2023research, torrejón2026teleoperation}.
    \item \textbf{Thiết bị phần cứng lỗi thời:} Việc sử dụng các thế hệ kính VR cũ (với độ phân giải thấp và tốc độ làm mới hạn chế) chưa thể đáp ứng được các yêu cầu khắt khe về nhận thức không gian và hạn chế sự tự do của người điều khiển \cite{hetrick2020comparing}.
\end{itemize}

\section{Đóng góp của luận văn và Bố cục}
\label{sec:contributions_outline}
Để giải quyết những thách thức nêu trên, luận văn này đề xuất phát triển một hệ thống điều khiển cánh tay robot từ xa hoàn chỉnh, ứng dụng những công nghệ phần cứng và phần mềm hiện đại nhất. Các đóng góp chính (Contributions) của luận văn được tóm tắt như sau:

\begin{itemize}
    \item \textbf{Xây dựng kiến trúc điều khiển thời gian thực dựa trên ROS2:} Luận văn khai thác sức mạnh của ROS2 (DDS middleware) cùng ROS2-Control để cải thiện tính thời gian thực, khắc phục rào cản về độ trễ của các hệ thống ROS1 cũ.
    \item \textbf{Phát triển ứng dụng VR nhập vai trên Meta Quest 3:} Tích hợp phần cứng hiện đại với Unity 6, xây dựng giao diện Digital Twin có khả năng đồng bộ trực tiếp chuyển động, cung cấp hình ảnh camera thời gian thực nhằm tối ưu hoá tải nhận thức cho người vận hành.
    \item \textbf{Thiết kế hệ thống Decoupled linh hoạt:} Đề xuất một kiến trúc tách biệt rõ ràng giữa tầng hiển thị/tương tác VR (Unity) và tầng điều khiển vật lý (ROS2/RC5 Controller), giúp hệ thống dễ dàng mở rộng và ứng dụng trên các loại cánh tay robot công nghiệp khác nhau.
\end{itemize}

\textbf{Bố cục của luận văn (Thesis Outline):} 
Chương 2 trình bày tổng quan các nghiên cứu liên quan về điều khiển robot qua VR và middleware. Chương 3 giới thiệu các cơ sở lý thuyết, bao gồm động học robot, cấu trúc ROS2, và công cụ Unity. Chương 4 đi sâu vào thiết kế kiến trúc phần mềm và hiện thực hệ thống điều khiển. Chương 5 trình bày và đánh giá các kết quả thực nghiệm đạt được. Cuối cùng, Chương 6 đưa ra kết luận và định hướng nghiên cứu trong tương lai.

% ---------------------------------------------------------

\chapter{CÁC NGHIÊN CỨU LIÊN QUAN}
\label{chap:related_works}

Trong những năm gần đây, sự giao thoa giữa robot học và công nghệ thực tế ảo đã thu hút được sự quan tâm lớn từ cộng đồng học thuật. Chương này sẽ phân tích các nghiên cứu liên quan, từ phương thức điều khiển, sự tích hợp giao diện VR, các giải pháp middleware truyền dẫn, và chỉ ra những khoảng trống nghiên cứu (Research Gaps) mà luận văn này hướng tới để giải quyết.

\section{Các phương pháp giao tiếp và điều khiển robot từ xa}
Phương pháp điều khiển robot từ xa truyền thống thường sử dụng Giao diện người dùng đồ họa (Graphical User Interface - GUI) kết hợp với chuột, bàn phím hoặc thiết bị ngoại vi chuyên dụng như SpaceMouse. Tuy nhiên, các giao diện này thường không phù hợp với các thao tác đòi hỏi độ chính xác không gian 3D. 

Trong nghiên cứu của Chen và cộng sự \cite{chen2023comparing}, tác giả đã thực hiện các bài đánh giá so sánh trực tiếp giữa GUI truyền thống, việc sử dụng Hand tracking (nhận diện cử chỉ tay) và VR Controller trên cánh tay robot JAKA Minicobo. Kết quả chỉ ra rằng VR Controller mang lại độ ổn định cao hơn, giảm hiện tượng nhiễu (jitter) tự nhiên của bàn tay và giúp người dùng hoàn thành nhiệm vụ nhanh chóng hơn. Dù vậy, việc lựa chọn phương thức điều khiển luôn đi kèm với bài toán về sự thoải mái công thái học và độ tin cậy của dữ liệu tín hiệu đầu vào.

\section{Tích hợp Thực tế ảo (VR) và Nhận thức tình huống (Situational Awareness)}
Nhận thức tình huống trong điều khiển từ xa đòi hỏi người vận hành phải thấu hiểu môi trường từ góc nhìn của robot. Để cung cấp một nhận thức không gian tốt hơn, nhiều nghiên cứu đã phát triển nền tảng Bản sao số (Digital Twin) bên trong VR \cite{torrejón2026teleoperation}.

Các nghiên cứu như của Torrejón và cộng sự \cite{torrejón2026teleoperation} hoặc hệ thống \textit{Vicarios} của Naceri và cộng sự \cite{naceri2021vicarios} cho thấy việc truyền luồng video (Video Streaming) và đám mây điểm (Point Cloud) từ không gian làm việc thực vào môi trường VR giúp người dùng dễ dàng phán đoán độ sâu và khoảng cách. Hơn nữa, việc cho phép người dùng thay đổi góc nhìn linh hoạt (teleporting/viewpoint changes) trong không gian VR giúp vượt qua những hạn chế do góc khuất vật lý (occlusion) mà các hệ thống camera tĩnh thường gặp phải \cite{naceri2021vicarios, su2022mixed}.

\section{Vấn đề Middleware và tối ưu hóa truyền dữ liệu}
Trung tâm của hệ thống Telerobotics nằm ở khả năng giao tiếp mượt mà giữa thiết bị VR và phần cứng Robot. ROS (Robot Operating System) từ lâu đã trở thành tiêu chuẩn trong công nghiệp và học thuật. Công cụ ROS-Sharp (ROS\#) được ứng dụng rộng rãi để kết nối engine Unity và ROS \cite{zhang2023research}. 

Tuy nhiên, khi truyền các dữ liệu lớn như hình ảnh hoặc point cloud, nền tảng ROS1 thường bộc lộ những hạn chế nghiêm trọng về độ trễ, đặc biệt là trong các hệ thống mạng Internet of Things (IoT) công nghiệp \cite{luu2025enhancing}. Các nghiên cứu gần đây đang có xu hướng dịch chuyển sang ROS2, sử dụng kiến trúc phân tán Data Distribution Service (DDS), nhằm tăng cường tính bảo mật và khả năng đáp ứng thời gian thực cho các tín hiệu điều khiển khắt khe \cite{singh2025comparative}.

\section{Khoảng trống nghiên cứu (Research Gaps)}
Qua việc phân tích các nghiên cứu kể trên, có thể rút ra một số khoảng trống công nghệ như sau:
\begin{enumerate}
    \item \textbf{Tính thời gian thực của Middleware phần mềm:} Mặc dù VR mang lại trải nghiệm ưu việt, phần lớn các hệ thống hiện tại \cite{hetrick2020comparing, chen2023comparing} vẫn được xây dựng trên ROS1, gây ra các giới hạn băng thông khi xử lý đa luồng dữ liệu (điều khiển + hình ảnh). Việc ứng dụng triệt để ROS2 cùng bộ công cụ ROS2-Control vào thực tế ảo để cải thiện độ trễ vẫn chưa được hiện thực hóa một cách phổ biến và toàn diện.
    \item \textbf{Giới hạn của thiết bị phần cứng VR:} Nhiều nghiên cứu trước đây bị giới hạn bởi các thiết bị thế hệ cũ (như HTC Vive đời đầu, hay Oculus Quest 2) \cite{hetrick2020comparing, chen2023comparing}. Các thiết bị này có hạn chế về chất lượng quang học hiển thị (Pass-through) và sức mạnh xử lý tính toán đồ hoạ cục bộ, làm giảm hiệu suất thiết lập Digital Twin với hình ảnh độ nét cao.
\end{enumerate}

Dựa trên những rào cản về công nghệ phần mềm (cơ chế Master/Slave của ROS1) và rào cản về phần cứng được rút ra từ các nghiên cứu trước, luận văn này đề xuất một kiến trúc hệ thống hiện đại hơn. Sự kết hợp giữa middleware ROS2, công cụ Unity 6, và thiết bị Meta Quest 3 hứa hẹn sẽ lấp đầy những khoảng trống này, tạo ra một trải nghiệm điều khiển thời gian thực có tính đồng bộ cao và trực quan nhất.
